\documentclass[a4paper,12pt,twoside]{memoir}

% Castellano
\usepackage[spanish,es-tabla]{babel}
\selectlanguage{spanish}
\usepackage[utf8]{inputenc}
\usepackage[T1]{fontenc}
\usepackage{lmodern} % Scalable font
\usepackage{microtype}
\usepackage{placeins}

\RequirePackage{booktabs}
\RequirePackage[table]{xcolor}
\RequirePackage{xtab}
\RequirePackage{multirow}

% Links
\usepackage[colorlinks]{hyperref}
\hypersetup{
	allcolors = {red}
}

% Ecuaciones
\usepackage{amsmath}

% Rutas de fichero / paquete
\newcommand{\ruta}[1]{{\sffamily #1}}

% Párrafos
\nonzeroparskip


% Imagenes
\usepackage{graphicx}
\newcommand{\imagen}[2]{
	\begin{figure}[!h]
		\centering
		\includegraphics[width=0.9\textwidth]{#1}
		\caption{#2}\label{fig:#1}
	\end{figure}
	\FloatBarrier
}

\newcommand{\imagenflotante}[2]{
	\begin{figure}%[!h]
		\centering
		\includegraphics[width=0.9\textwidth]{#1}
		\caption{#2}\label{fig:#1}
	\end{figure}
}



% El comando \figura nos permite insertar figuras comodamente, y utilizando
% siempre el mismo formato. Los parametros son:
% 1 -> Porcentaje del ancho de página que ocupará la figura (de 0 a 1)
% 2 --> Fichero de la imagen
% 3 --> Texto a pie de imagen
% 4 --> Etiqueta (label) para referencias
% 5 --> Opciones que queramos pasarle al \includegraphics
% 6 --> Opciones de posicionamiento a pasarle a \begin{figure}
\newcommand{\figuraConPosicion}[6]{%
  \setlength{\anchoFloat}{#1\textwidth}%
  \addtolength{\anchoFloat}{-4\fboxsep}%
  \setlength{\anchoFigura}{\anchoFloat}%
  \begin{figure}[#6]
    \begin{center}%
      \Ovalbox{%
        \begin{minipage}{\anchoFloat}%
          \begin{center}%
            \includegraphics[width=\anchoFigura,#5]{#2}%
            \caption{#3}%
            \label{#4}%
          \end{center}%
        \end{minipage}
      }%
    \end{center}%
  \end{figure}%
}

%
% Comando para incluir imágenes en formato apaisado (sin marco).
\newcommand{\figuraApaisadaSinMarco}[5]{%
  \begin{figure}%
    \begin{center}%
    \includegraphics[angle=90,height=#1\textheight,#5]{#2}%
    \caption{#3}%
    \label{#4}%
    \end{center}%
  \end{figure}%
}
% Para las tablas
\newcommand{\otoprule}{\midrule [\heavyrulewidth]}
%
% Nuevo comando para tablas pequeñas (menos de una página).
\newcommand{\tablaSmall}[5]{%
 \begin{table}
  \begin{center}
   \rowcolors {2}{gray!35}{}
   \begin{tabular}{#2}
    \toprule
    #4
    \otoprule
    #5
    \bottomrule
   \end{tabular}
   \caption{#1}
   \label{tabla:#3}
  \end{center}
 \end{table}
}

%
% Nuevo comando para tablas pequeñas (menos de una página).
\newcommand{\tablaSmallSinColores}[5]{%
 \begin{table}[H]
  \begin{center}
   \begin{tabular}{#2}
    \toprule
    #4
    \otoprule
    #5
    \bottomrule
   \end{tabular}
   \caption{#1}
   \label{tabla:#3}
  \end{center}
 \end{table}
}

\newcommand{\tablaApaisadaSmall}[5]{%
\begin{landscape}
  \begin{table}
   \begin{center}
    \rowcolors {2}{gray!35}{}
    \begin{tabular}{#2}
     \toprule
     #4
     \otoprule
     #5
     \bottomrule
    \end{tabular}
    \caption{#1}
    \label{tabla:#3}
   \end{center}
  \end{table}
\end{landscape}
}

%
% Nuevo comando para tablas grandes con cabecera y filas alternas coloreadas en gris.
\newcommand{\tabla}[6]{%
  \begin{center}
    \tablefirsthead{
      \toprule
      #5
      \otoprule
    }
    \tablehead{
      \multicolumn{#3}{l}{\small\sl continúa desde la página anterior}\\
      \toprule
      #5
      \otoprule
    }
    \tabletail{
      \hline
      \multicolumn{#3}{r}{\small\sl continúa en la página siguiente}\\
    }
    \tablelasttail{
      \hline
    }
    \bottomcaption{#1}
    \rowcolors {2}{gray!35}{}
    \begin{xtabular}{#2}
      #6
      \bottomrule
    \end{xtabular}
    \label{tabla:#4}
  \end{center}
}

%
% Nuevo comando para tablas grandes con cabecera.
\newcommand{\tablaSinColores}[6]{%
  \begin{center}
    \tablefirsthead{
      \toprule
      #5
      \otoprule
    }
    \tablehead{
      \multicolumn{#3}{l}{\small\sl continúa desde la página anterior}\\
      \toprule
      #5
      \otoprule
    }
    \tabletail{
      \hline
      \multicolumn{#3}{r}{\small\sl continúa en la página siguiente}\\
    }
    \tablelasttail{
      \hline
    }
    \bottomcaption{#1}
    \begin{xtabular}{#2}
      #6
      \bottomrule
    \end{xtabular}
    \label{tabla:#4}
  \end{center}
}

%
% Nuevo comando para tablas grandes sin cabecera.
\newcommand{\tablaSinCabecera}[5]{%
  \begin{center}
    \tablefirsthead{
      \toprule
    }
    \tablehead{
      \multicolumn{#3}{l}{\small\sl continúa desde la página anterior}\\
      \hline
    }
    \tabletail{
      \hline
      \multicolumn{#3}{r}{\small\sl continúa en la página siguiente}\\
    }
    \tablelasttail{
      \hline
    }
    \bottomcaption{#1}
  \begin{xtabular}{#2}
    #5
   \bottomrule
  \end{xtabular}
  \label{tabla:#4}
  \end{center}
}



\definecolor{cgoLight}{HTML}{EEEEEE}
\definecolor{cgoExtralight}{HTML}{FFFFFF}

%
% Nuevo comando para tablas grandes sin cabecera.
\newcommand{\tablaSinCabeceraConBandas}[5]{%
  \begin{center}
    \tablefirsthead{
      \toprule
    }
    \tablehead{
      \multicolumn{#3}{l}{\small\sl continúa desde la página anterior}\\
      \hline
    }
    \tabletail{
      \hline
      \multicolumn{#3}{r}{\small\sl continúa en la página siguiente}\\
    }
    \tablelasttail{
      \hline
    }
    \bottomcaption{#1}
    \rowcolors[]{1}{cgoExtralight}{cgoLight}

  \begin{xtabular}{#2}
    #5
   \bottomrule
  \end{xtabular}
  \label{tabla:#4}
  \end{center}
}


















\graphicspath{ {./img/} }

% Capítulos
\chapterstyle{bianchi}
\newcommand{\capitulo}[2]{
	\setcounter{chapter}{#1}
	\setcounter{section}{0}
	\chapter*{#2}
	\addcontentsline{toc}{chapter}{#2}
	\markboth{#2}{#2}
}

% Apéndices
\renewcommand{\appendixname}{Apéndice}
\renewcommand*\cftappendixname{\appendixname}

\newcommand{\apendice}[1]{
	%\renewcommand{\thechapter}{A}
	\chapter{#1}
}

\renewcommand*\cftappendixname{\appendixname\ }

% Formato de portada
\makeatletter
\usepackage{xcolor}
\newcommand{\tutor}[1]{\def\@tutor{#1}}
\newcommand{\course}[1]{\def\@course{#1}}
\definecolor{cpardoBox}{HTML}{E6E6FF}
\def\maketitle{
  \null
  \thispagestyle{empty}
  % Cabecera ----------------
\noindent\includegraphics[width=\textwidth]{cabecera}\vspace{1cm}%
  \vfill
  % Título proyecto y escudo informática ----------------
  \colorbox{cpardoBox}{%
    \begin{minipage}{.8\textwidth}
      \vspace{.5cm}\Large
      \begin{center}
      \textbf{TFG del Grado en Ingeniería Informática}\vspace{.6cm}\\
      \textbf{\LARGE\@title{}}
      \end{center}
      \vspace{.2cm}
    \end{minipage}

  }%
  \hfill\begin{minipage}{.20\textwidth}
    \includegraphics[width=\textwidth]{escudoInfor}
  \end{minipage}
  \vfill
  % Datos de alumno, curso y tutores ------------------
  \begin{center}%
  {%
    \noindent\LARGE
    Presentado por \@author{}\\ 
    en Universidad de Burgos --- \@date{}\\
    Tutor: \@tutor{}\\
  }%
  \end{center}%
  \null
  \cleardoublepage
  }
\makeatother

\newcommand{\nombre}{Nombre del alumno} %%% cambio de comando

% Datos de portada
\title{título del TFG}
\author{\nombre}
\tutor{nombre tutor}
\date{\today}

\begin{document}

\maketitle


\newpage\null\thispagestyle{empty}\newpage


%%%%%%%%%%%%%%%%%%%%%%%%%%%%%%%%%%%%%%%%%%%%%%%%%%%%%%%%%%%%%%%%%%%%%%%%%%%%%%%%%%%%%%%%
\thispagestyle{empty}


\noindent\includegraphics[width=\textwidth]{cabecera}\vspace{1cm}

\noindent D. nombre tutor, profesor del departamento de nombre departamento, área de nombre área.

\noindent Expone:

\noindent Que el alumno D. \nombre, con DNI dni, ha realizado el Trabajo final de Grado en Ingeniería Informática titulado título de TFG. 

\noindent Y que dicho trabajo ha sido realizado por el alumno bajo la dirección del que suscribe, en virtud de lo cual se autoriza su presentación y defensa.

\begin{center} %\large
En Burgos, {\large \today}
\end{center}

\vfill\vfill\vfill

% Author and supervisor
\begin{minipage}{0.45\textwidth}
\begin{flushleft} %\large
Vº. Bº. del Tutor:\\[2cm]
D. nombre tutor
\end{flushleft}
\end{minipage}
\hfill
\begin{minipage}{0.45\textwidth}
\begin{flushleft} %\large
Vº. Bº. del co-tutor:\\[2cm]
D. nombre co-tutor
\end{flushleft}
\end{minipage}
\hfill

\vfill

% para casos con solo un tutor comentar lo anterior
% y descomentar lo siguiente
%Vº. Bº. del Tutor:\\[2cm]
%D. nombre tutor


\newpage\null\thispagestyle{empty}\newpage




\frontmatter

% Abstract en castellano
\renewcommand*\abstractname{Resumen}
\begin{abstract}
En este primer apartado se hace una \textbf{breve} presentación del tema que se aborda en el proyecto.
\end{abstract}

\renewcommand*\abstractname{Descriptores}
\begin{abstract}
Palabras separadas por comas que identifiquen el contenido del proyecto Ej: servidor web, buscador de vuelos, android \ldots
\end{abstract}

\clearpage

% Abstract en inglés
\renewcommand*\abstractname{Abstract}
\begin{abstract}
A \textbf{brief} presentation of the topic addressed in the project.
\end{abstract}

\renewcommand*\abstractname{Keywords}
\begin{abstract}
keywords separated by commas.
\end{abstract}

\clearpage

% Indices
\tableofcontents

\clearpage

\listoffigures

\clearpage

\listoftables
\clearpage

\mainmatter
\capitulo{1}{Introducción}

Descripción del contenido del trabajo y del estrucutra de la memoria y del resto de materiales entregados.

\capitulo{2}{Objetivos del proyecto}

Con el fin de solucionar los problemas concernientes a la docencia antes mencionados, comenzaremos por definir los resultados esperados. Necesitaremos un visor que represente gráficamente los modelos así como un sistema que sea capaz de gestionar la información que se añadirá sobre los mismos. Además, tendremos que limitar el acceso a la aplicación al público esperado mediante un sistema de validación de usuarios. Esta restricción se da para mantener los modelos virtuales seguros, de forma similar a los físicos; solamente podrán acceder a ellos los alumnos autorizados, así como profesores e investigadores a los que se desee dar acceso.

Durante la ejecución del proyecto, añadiremos como objetivo la realización de análisis sobre archivos de los modelos. Con esto, lograremos descubrir aquellos que hayan sido generados de una manera poco correcta (con puntos en el fichero que no formen parte del modelo).

\section{Visualización}
Lo primero que necesitamos es un sistema que permita visualizar los modelos\footnote{En formato PLY (avanzar a sección~\ref{sec:formato-ply}), poseen gran cantidad de detalles y son muy realistas} que el profesor provea, y queremos que éstos no se compartan fácilmente. De esta manera, necesitamos que se carguen de forma transparente al usuario, que éste no tenga en la medida de lo posible acceso a los ficheros que lo contienen. Teniendo ambas cosas en cuenta, emplearemos algún tipo de \textit{framework} que nos permita la visualización, interacción y carga de los modelos en el formato que tenemos. Un punto importante para su uso reside en evitar la necesidad de instalación, así que pensamos en realizar el desarrollo de una página \textit{web} con el fin de tener una gran portabilidad; con un navegador y conexión a Internet los usuarios pueden comenzar.

\section{Servidor}
Con el fin de dar servicio sobre los diferentes contenidos, necesitamos un sistema que gestione los elementos de que la aplicación se compone. Por ello, crearemos un servidor basado en el \textit{microframework} Flask que nos permita generar de forma dinámica los contenidos.

\section{Control de acceso}
Dado que la aplicación está destinada a alumnos del Grado en Historia y Patrimonio, además de necesitar que los modelos sean lo más privados posible, es imperativo que la aplicación tenga algún sistema de evitar accesos no autorizados. Con este fin, crearemos un control de acceso mediante \textit{login}.

Éste restringirá el uso de nuestra herramienta verificando que, en efecto, el usuario pertenezca a la comunidad universitaria de la UBU mediante una llamada a la \textit{API} de UBUVirtual, que está basada en Moodle\footnote{Tal y como aparece mencionado en \url{https://ubuvirtual.ubu.es/mod/page/view.php?id=18&lang=es}}~\cite{github:moodle}. Adicionalmente, se verificará que pertenezca a una lista de usuarios autorizados.

\section{Gestión de la información añadida}
Además del modelo, necesitamos representar la información que generaremos para los mismos. Proveeremos de los medios adecuados para representar gráficamente y realizar la persistencia de dichos elementos, que serán anotaciones y medidas.

Gráficamente, se podrán ver junto al modelo generado en forma de esferas representando los puntos donde están colocadas en el caso de las anotaciones. En el caso de las medidas, serán esferas para los extremos junto una línea que los una.

Con el fin de que esta información se pueda intercambiar fácilmente, se diseñará una estructura de datos a seguir para poder almacenarlos en forma de archivo. En nuestro caso, diseñaremos una definición de \textit{JSON} para dicha necesidad.


\section{Análisis de los modelos}\label{sec:analisis-modelos}
En el proceso de generación del modelo\footnote{Los ficheros \textit{PLY} se consiguen a partir de múltiples fotografías obtenidas con un escáner 3D, el proceso automático no es perfecto y necesita que los investigadores reparen el modelo 3D obtenido. Incluso tras la revisión manual, en muchas ocasiones siguen apareciendo puntos que no se corresponden con el modelo}, es posible que haya pequeños errores, puntos que son ruido porque se capturan por error y no pertenecen al mismo. Aunque aparecen dichos puntos en el fichero del modelo, no pertenecen a ninguno de polígonos con los que se representa la figura, son datos sin valor alguno. Si no se veían, ¿porqué entonces nos íbamos a dar cuenta de su presencia? Con estos datos, la forma estándar de calcular las dimensiones del modelo eran imprecisas, y nos basábamos en dichas dimensiones para centrar adecuadamente el modelo durante la fase de carga en pantalla.

Con el fin de conseguir una visualización más acorde al ideal, nos vimos en la necesidad de descubrir qué puntos eran los que nos molestaban para librarnos de ellos posteriormente. Dichas herramientas se pueden encontrar en el directorio \texttt{<<Resources>>} del proyecto. En el anexo <<Manual del programador>> se explica con detalle el uso de las mismas.

Nótese que esta es una herramienta suplementaria a eliminar el ruido de fondo de forma manual, como aparece documentado en~\cite{wiki:removing-background-noise}.
\capitulo{3}{Conceptos teóricos}

En aquellos proyectos que necesiten para su comprensión y desarrollo de unos conceptos teóricos de una determinada materia o de un determinado dominio de conocimiento, debe existir un apartado que sintetice dichos conceptos.

Algunos conceptos teóricos de \LaTeX \footnote{Créditos a los proyectos de Álvaro López Cantero: Configurador de Presupuestos y Roberto Izquierdo Amo: PLQuiz}.

\section{Secciones}

Las secciones se incluyen con el comando section.

\subsection{Subsecciones}

Además de secciones tenemos subsecciones.

\subsubsection{Subsubsecciones}

Y subsecciones. 


\section{Referencias}

Las referencias se incluyen en el texto usando cite \cite{wiki:latex}. Para citar webs, artículos o libros \cite{koza92}.


\section{Imágenes}

Se pueden incluir imágenes con los comandos standard de \LaTeX, pero esta plantilla dispone de comandos propios como por ejemplo el siguiente:

\imagen{escudoInfor}{Autómata para una expresión vacía}



\section{Listas de items}

Existen tres posibilidades:

\begin{itemize}
	\item primer item.
	\item segundo item.
\end{itemize}

\begin{enumerate}
	\item primer item.
	\item segundo item.
\end{enumerate}

\begin{description}
	\item[Primer item] más información sobre el primer item.
	\item[Segundo item] más información sobre el segundo item.
\end{description}
	
\begin{itemize}
\item 
\end{itemize}

\section{Tablas}

Igualmente se pueden usar los comandos específicos de \LaTeX o bien usar alguno de los comandos de la plantilla.

\tablaSmall{Herramientas y tecnologías utilizadas en cada parte del proyecto}{l c c c c}{herramientasportipodeuso}
{ \multicolumn{1}{l}{Herramientas} & App AngularJS & API REST & BD & Memoria \\}{ 
HTML5 & X & & &\\
CSS3 & X & & &\\
BOOTSTRAP & X & & &\\
JavaScript & X & & &\\
AngularJS & X & & &\\
Bower & X & & &\\
PHP & & X & &\\
Karma + Jasmine & X & & &\\
Slim framework & & X & &\\
Idiorm & & X & &\\
Composer & & X & &\\
JSON & X & X & &\\
PhpStorm & X & X & &\\
MySQL & & & X &\\
PhpMyAdmin & & & X &\\
Git + BitBucket & X & X & X & X\\
Mik\TeX{} & & & & X\\
\TeX{}Maker & & & & X\\
Astah & & & & X\\
Balsamiq Mockups & X & & &\\
VersionOne & X & X & X & X\\
} 

\capitulo{4}{Técnicas y herramientas}

En las siguientes secciones vamos a documentar algunos elementos con los que justificar las decisiones tomadas sobre las librerías y herramientas de trabajo empleadas.

\section{Comparativa de \textit{IDE}}
Para facilitar el desarrollo del \textit{plugin} de \textit{JavaScript}, queremos un \textit{IDE} con función de autocompletar. Aunque varios disponen de dicha función para \textit{JavaScript} como lenguaje, pocos son capaces de aportar esta característica a las bibliotecas que queremos usar. Entre estos pocos se encuentran Sublimetext\footnote{\url{https://www.sublimetext.com/}}, que lo soporta mediante una extensión de terceros\footnote{\url{https://github.com/blackjk3/threejs-sublime}}, aunque éste es un \textit{IDE} con versión de prueba y finalmente tendríamos que pagar una licencia. Otro que permite autocompletar en bibliotecas (en nuestro caso, three.js\footnote{\url{https://threejs.org/}}) es Visual Studio Code. Para habilitar esta característica  tendríamos que seguir una pequeña guía\footnote{\url{https://code.visualstudio.com/docs/languages/javascript\#_typescript-definition-files-typings}}. Otro gran \textit{IDE} como LightTable lo hemos descartado por no tener autocompletar para three.js.

\section{Editor de documentación \LaTeX}
La documentación de la memoria y anexos la vamos a realizar basándonos en la plantilla que nos provee la UBU\footnote{\url{https://github.com/ubutfgm/plantillaLatex}}.

De los diferentes editores de \LaTeX que existen disponibles (consideramos TeXstudio y Texmaker), nos hemos decantado por el primero al resultarnos más fácil de usar la interfaz.

\section{Librería 3D en \textit{JavaScript}: three.js}
Al inicio del proyecto, se buscaba una librería que permitiese la visualización de elementos tridimensionales en un navegador \textit{web}. Deseábamos que nuestro visor fuese capaz de cargar archivos con extensión \textit{STL}, por lo que valoramos emplear babylon.js\footnote{\url{http://www.babylonjs.com/}}, openjscad\footnote{\url{https://openjscad.org/}} y {three.js\footnote{\url{https://threejs.org/}}.

Aunque todas ellas tenían dicho soporte, también valoramos otras características, como el número de proyectos hechos a partir de ellas, la cantidad de documentación disponible para las mismas, o los módulos disponibles. Como consecuencia, nos decantamos finalmente por three.js.

Un tiempo después de iniciar el proyecto, se cambió un requisito y desde entonces vamos a cargar archivos de tipo \textit{PLY} y no \textit{STL}. Es por esto que fue un acierto escoger three.js: soportar una gran cantidad de formatos hizo poco traumático el nuevo requisito. Con otras opciones probablemente tendríamos que haber empezado de cero en otro \textit{framework}, o si acaso crear un cargador que fuese compatible.

\section{Librería GUI en \textit{JavaScript}: dat.GUI}
Conocimos esta librería por ser prácticamente omnipresente en una gran cantidad de ejemplos y tutoriales de three.js. Existen algunas páginas\footnote{\url{http://workshop.chromeexperiments.com/examples/gui/}}\footnote{\url{https://davidwalsh.name/dat-gui}} en las cuales podemos comprender perfectamente la utilidad así como la manera de configurar dichos controles, aunque echamos en falta documentación de algún tipo en el propio proyecto.

Aunque existe un mayor número de elementos, para nuestro proyecto hemos empleado: carpetas, para mantener organizados los elementos de diferentes categorías, \textit{checkbox} o casilla de confirmación para activar y desactivar el asistente de malla, y \textit{dropdown list} o lista desplegable para mostrar diferentes escalas para escoger.

Su simplicidad radica en no tener que configurar los elementos para que sean un desplegable, una casilla de confirmación o incluso un cuadro de texto, sino que según el tipo de parámetro que se le pase, el elemento que genera en el menú se elige de forma transparente al programador.

Sin embargo, la facilidad para configurar los elementos de la que hace gala queda empañada por la escasa manipulación que permite en los diferentes elementos que pueden usarse: no podemos integrar varios botones en un mismo elemento, ni cambiar su estilo. Tampoco podemos incluir elementos que simplemente sean seleccionables y no botones o campos de texto, aquí no tenemos capacidad de elección.

Finalmente, esto no quiere decir que no sirva para lo que ha sido concebida: tener un menú que permita modificaciones de parámetros a la par que dicho menú sea de fácil desarrollo.

\section{Librería funcionalidades en \textit{JavaScript}: jQuery}
jQuery\footnote{\url{http://jquery.com/}} es una librería \textit{JavaScript} que permite manipular, manejar eventos y otras muchas funcionalidades mediante una \textit{API} simple de usar al mismo tiempo que mantiene una gran compatibilidad entre navegadores \textit{web}.

<<\$>>. Es el símbolo que define esta librería. A través de éste, llamamos a la función principal que nos permite realizar selecciones de diferentes elementos del \textit{DOM} que queramos manejar, y todo ello empleando un estándar conocido por los desarrolladores \textit{web}, como son los selectores \textit{CSS}.
Solamente pasando una cadena con un selector, podemos realizar diferentes acciones que, de otra manera hubiesen sido tediosas de realizar en \textit{JavaScript} puro. Por ejemplo, en nuestra aplicación, para deseleccionar todos los elementos de una lista en un menú, hacemos lo siguiente:

\verb|$("#menu li.ui-selected").removeClass("ui-selected");|

Lo que hacemos es seleccionar todos los elementos \texttt{<<li>>} que se encuentren seleccionados dentro de \texttt{<<\#menu>>} y deseleccionarlos.

\section{Librería \textit{GUI} en \textit{JavaScript}: jQueryUI}
jQueryUI\footnote{\url{http://jqueryui.com/}}, al igual que jQuery, está mantenida por la jQuery Foundation\footnote{\url{https://jquery.org/}}. Se basa en la sintaxis de su hermana (depende de ella) para aportar funcionalidades sobre la parte gráfica de las páginas \textit{web}. Así, permite crear listas, botones personalizados, elementos seleccionables y ventanas emergentes, además de otros elementos que no hemos empleado en nuestra aplicación.

Y no es solamente que nos permita crear estos elementos gráficos, sino que además nos proporciona una interfaz para enlazar los eventos que generan con las acciones que nosotros queramos.

Gracias a estas características la escogimos como librería para permitirnos un desarrollo rápido y eficiente para los elementos de la interfaz en el visor.

\section{Diagramas \textit{UML}: draw.io}
Para editar los diagramas \textit{UML} de clases, empleamos la herramienta en línea draw.io\footnote{\url{https://www.draw.io/}}, que nos provee suficiente funcionalidad sin necesidad de instalar ninguna herramienta extra.

\section{Documentar \textit{JSON}}
Con el fin de mostrar el esquema de los archivos \textit{JSON} que se iban a importar, nos vemos en la necesidad de destacar que la plantilla empleada para resaltar \textit{JSON} ha sido obtenida de \TeX\space\textit{stackexchange}\footnote{\url{https://tex.stackexchange.com/questions/83085/how-to-improve-listings-display-of-json-files}}.

\section{Metodología Scrum}
Emplearemos Scrum durante el proyecto apoyándonos en las características propias de GitHub, así como el añadido que genera ZenHub\footnote{\url{https://www.zenhub.com/}} sobre él.
\capitulo{5}{Aspectos relevantes del desarrollo del proyecto}

Este apartado pretende recoger los aspectos más interesantes del desarrollo del proyecto, comentados por los autores del mismo.
Debe incluir desde la exposición del ciclo de vida utilizado, hasta los detalles de mayor relevancia de las fases de análisis, diseño e implementación.
Se busca que no sea una mera operación de copiar y pegar diagramas y extractos del código fuente, sino que realmente se justifiquen los caminos de solución que se han tomado, especialmente aquellos que no sean triviales.
Puede ser el lugar más adecuado para documentar los aspectos más interesantes del diseño y de la implementación, con un mayor hincapié en aspectos tales como el tipo de arquitectura elegido, los índices de las tablas de la base de datos, normalización y desnormalización, distribución en ficheros3, reglas de negocio dentro de las bases de datos (EDVHV GH GDWRV DFWLYDV), aspectos de desarrollo relacionados con el WWW...
Este apartado, debe convertirse en el resumen de la experiencia práctica del proyecto, y por sí mismo justifica que la memoria se convierta en un documento útil, fuente de referencia para los autores, los tutores y futuros alumnos.

\capitulo{6}{Trabajos relacionados}

Para conseguir un correcto desarrollo de nuestro proyecto, observamos algunas de las herramientas más representativas que hemos encontrado en el mercado y las comparamos, tanto entre ellas como con la nuestra.

\section{Comparativa de visores \textit{STL}}
Haremos una comparativa de algunos de los visores \textit{STL} más relevantes. Casi todos los que vamos a analizar son gratuitos y funcionan de forma \textit{online}, exceptuando algunos. Aunque estos visores normalmente se emplean \textit{online}, es posible utilizarlos de forma local en el navegador incluyéndolos en un archivo \textit{html}.

Ponemos un resumen en la tabla \ref{tabla:caracteristicas-herramientas}.

\tablaSmall
{Resumen de características}
{lccccc}
{caracteristicas-herramientas}
{Nombre herramienta & STL & PLY & Plataforma & Código abierto & Activo \\}
{
	ViewSTL & Si & No & WebGL & No (pago) & Si \\
	OpenJSCAD & Si & No & WebGL & Si & Si \\ 
	Mesh Viewer & No & Si & Escritorio & Si  & No \\ 
	Open 3D Model Viewer & Si & Si & Escritorio & Si & Si \\ 
	Pointcloud-PLY-Viewer & No & Si & Flash & Si & No \\ 
	Online 3D PLY & No & Si & WebGL & No (pago) & Si \\ 
	Autodesk A360 Viewer & Si & Si & Plug-in web & No (pago) & Si \\ 
	Nuestro visor & No & Si & WebGL & Si & Si \\ 
}

\subsection{ViewSTL}
Es muy simple (el más simple de todos los que hemos examinado)\footnote{\url{http://www.viewstl.com/}}, solamente permite la visualización de archivos en formato \textit{STL} (escondido en el \textit{FAQ}\footnote{\url{http://www.viewstl.com/faq/}} dice que soportan \textit{OBJ}). También es el único que no realiza actividad de red para cargar el archivo, por lo que cumple con lo que promete y no sube nuestros archivos a la red.
Sin embargo, tiene algunas opciones no observadas en los rivales, como definir las unidades de las dimensiones del archivo. Ya que el formato \textit{STL} no contiene las unidades en las que se definen las coordenadas, al menos podemos añadir la información a posteriori.
Otra característica relevante es la de definir la orientación del modelo desde un menú desplegable. Puede que no todos los modelos vengan capturados desde el mismo ángulo y, teniendo en cuenta las limitaciones que hay respecto a la rotación de las cámaras, es posible que nos sea de utilidad.
También podemos elegir el color del modelo, ya que aunque \textit{STL} soporte definir color, no es en absoluto estándar~\cite{wiki:color-in-STL} y normalmente deberemos definirlo nosotros.
La última característica que vamos a mencionar es la opción de tomar una instantánea de lo que estemos visualizando en ese momento podría ser útil a la hora de realizar imágenes para diapositivas o similares.

\subsection{Autodesk A360 Viewer}
Creado por Autodesk, el visor\footnote{\url{https://a360.autodesk.com/viewer}} posee un aspecto profesional y de escritorio.
A destacar el asistente de medida, bastante visual e intuitivo y el asistente de perspectiva, que ayuda a rotar la pieza de la manera deseada.
Permite elegir entre diferentes maneras de rotar (órbitas), cambiar parámetros de la cámara, añadir planos que corten la pieza para ver el corte en sección, que resalta entre las demás por ser de nuestro interés.
Finalmente, podemos señalar que cambia el icono del puntero según la actividad que estemos realizando (rotando, \textit{zoom}, desplazando, midiendo\dots).

\subsection{OpenJSCAD}
Es el visor que nos permite ahondar más en los detalles del modelo. Al importar un archivo \textit{STL}, pasa primero por un asistente que transforma un archivo \textit{STL} (tanto binario como \textit{ASCII}) en el conjunto de coordenadas y triángulos de los que está compuesto reorganizando pues los datos de una manera que el propio OpenJSCAD\footnote{\url{http://openjscad.org/}} comprende. Tras esta <<interpretación>> del objeto, tenemos a nuestra disposición una versión editable de los datos en crudo, además de poder visualizar cambios sobre los mismos en tiempo real.
Como último punto, permite realizar la exportación de la figura visualizada en diferentes formatos.


\section {Comparativa de visores \textit{PLY}}
Tras la modificación del requisito del tipo de archivos que nuestro programa debe aceptar de \textit{STL} a \textit{PLY}~\cite{wiki:PLY}, necesitamos examinar los visores de archivos \textit{PLY} disponibles.

Nos hallamos desconcertados ante el aparente poco interés de la comunidad de código abierto respecto a este formato, puesto que no encontramos complementos para emplear en el navegador de forma gratuita, aunque vemos alguna excepción. La mayor parte de los programas que existen para el navegador son servicios de pago, y los programas de escritorio no son multiplataforma. Terminamos el análisis descorazonados por la inexistente oferta en el mercado, además de no poder comprobar de primera mano el funcionamiento de los mismos; o son de pago, o muy complicados de instalar. Por ello, exponemos un breve repaso a cada uno de ellos.

\subsection{Pointcloud-PLY-Viewer}
Ésta es la única implementación para navegador y de código abierto que hemos conseguido encontrar\footnote{\url{https://github.com/ktuite/Pointcloud-PLY-viewer}}, pero desafortunadamente se basa en el ya vetusto Adobe Flash y no ha sido actualizado recientemente (4 años).

\subsection{Mesh Viewer}
Es una aplicación de escritorio\footnote{\url{http://mview.sourceforge.net/}} multiplataforma (Linux, Windows), aunque solamente aparece el código fuente y no hemos conseguido obtener una compilación funcional. La compatibilidad con \textit{PLY} es limitada, pues solamente soporta su versión \textit{ASCII}.

\subsection{Open 3D Model Viewer}
Aplicación de escritorio\footnote{\url{http://www.open3mod.com/}} que únicamente está disponible para Windows con instalador y con versión portable. En principio permite mucho más que solamente visualizaciones, como realizar animaciones con los modelos, además de edición compleja. Tampoco lo hemos podido probar por no poseer licencia de Windows.

\subsection{Online 3D PLY Viewer}
Es una herramienta\footnote{\url{https://www.afanche.com/online-3d-ply-viewer}} que, como su propio nombre indica, permite la visualización \textit{online}. Sin embargo, hay que comprar un \textit{token} a la empresa propietaria para el uso de la misma.

Aunque les preguntamos el día 2 de marzo de 2017 acerca del precio, a día \today \space aún no hemos conseguido respuesta.

\section{African Fossils}
Un trabajo similar al nuestro sería el de African Fossils\footnote{\url{http://africanfossils.org/}}, aunque con un enfoque diferente.

Coincide en mostrar de forma tridimensional modelos, y se apoya en un plugin de Autodesk. Salvando las diferencias entre los tipos de modelos, añade a éstos explicaciones en texto, siendo éstos los únicos elementos que nos proporcionan información sobre el modelo. Permite descargar los modelos para imprimirlos en 3D, y también tienen disponible funciones de redes sociales para dar visibilidad al proyecto.


\capitulo{7}{Conclusiones y Líneas de trabajo futuras}

En esta parte expondremos las conclusiones extraídas a lo largo del trabajo, así como posibles líneas de trabajo futuro con las que continuar el proyecto.

\section{Conclusiones}
Con la ejecución del proyecto, hemos podido extraer una serie de conclusiones:
\begin{itemize}
	\item Hemos conseguido cumplir los objetivos iniciales del proyecto, además del añadido para realizar el análisis de los modelos. Ver sección~\ref{sec:analisis-modelos}.
	\item Que gracias a los conocimientos adquiridos durante la titulación, hemos creado una estructura para el código que facilitará los cambios que se realicen en un futuro.
	\item El empleo de tecnologías y herramientas de los cuales no teníamos experiencia ha retrasado el progreso, pero serán de utilidad a futuro (en especial \LaTeX).
	\item Precisamente el desconocimiento del área sobre el que versa el trabajo ha hecho bastante complicada la estimación del tiempo de las tareas, aunque el uso de metodologías ágiles nos ha aportado la flexibilidad suficiente para abordar dichos imprevistos.
\end{itemize}



\section{Lineas de trabajo futuras}
Ahora pasaremos a explicar algunos aspectos en los que el proyecto puede evolucionar en un futuro. Algunos de ellos se ven reflejados en el apartado \texttt{<<Issues>>} de GitHub.

\subsection{Despliegue en servidor de la UBU}
Sería una buena idea integrar de forma permanente el servicio en un servidor de la universidad por varios motivos. El primero sería la capacidad de mantenerlo como cualquier otro servicio propio de la universidad. Otro sería la seguridad de los datos: al ser los modelos valiosos para la institución, la mejor manera para mantener la privacidad de los mismos es mediante el alojamiento propio.

\subsection{Integración con Moodle}
Otro nivel de integración sería con los servicios ofrecidos actualmente en UBUVirtual. Uno de ellos sería el sistema de cuestionarios que posee. De esta manera, los usuarios podrán realizar los ejercicios en un menor número de pasos y concentrarse mejor en el ejercicio en sí mismo.

\subsection{Perfiles de usuario}
En el momento presente, no existe ningún tipo de perfil de usuario: todos tienen los mismos permisos. Por ello, será necesario implementar algún tipo de control a dicho nivel, para que los profesores o aquellos que administren el contenido sean capaces de limitar y conceder acceso a diferentes elementos de la aplicación.

\subsection{Comercialización}
Finalmente podemos mencionar la posibilidad de monetizar la aplicación, teniendo además varias posibilidades.

Una de ellas sería la de ofrecer acceso a los modelos a otros investigadores o profesores en otras instituciones, probablemente mediante un modelo de suscripción. De esta manera se podría facilitar el trabajo colaborativo sobre el conjunto de modelos.

Y la otra idea es la de desplegar nuestro proyecto sobre otras instituciones con similares necesidades pudiendo añadir además nuevas características, como el soporte de nuevos formatos de archivo, visualizar elementos diferentes como moléculas o nanomateriales\footnote{\url{http://iesbinef.educa.aragon.es/fiqui/jmol/organica.htm?_USE=HTML5}}, etc.


\bibliographystyle{plain}
\bibliography{bibliografia}

\end{document}